\documentclass[11pt,a4paper]{article}
\usepackage{fontspec}
\usepackage{polyglossia}
\setmainlanguage{polish}
\setmainfont{Latin Modern Roman}
\setsansfont{Latin Modern Sans}
\setmonofont{Latin Modern Mono}
\usepackage{geometry}
\usepackage{graphicx}
\usepackage{booktabs}
\usepackage{array}
\usepackage{float}
\usepackage{caption}
\usepackage{hyperref}
\usepackage{xcolor}
\usepackage{listings}
\lstdefinestyle{pythonstyle}{
    language=Python,
    basicstyle=\small\ttfamily,
    breaklines=true,
    columns=fullflexible,
    keepspaces=true,
    frame=single,
    rulecolor=\color{black},
    showstringspaces=false
}
\lstset{style=pythonstyle}
\geometry{margin=2.2cm}
\setlength{\parindent}{0pt}
\setlength{\parskip}{0.55em}
\renewcommand{\arraystretch}{1.15}
\sloppy

\title{Laboratorium 6\\Dane sonarowe BSS i klasyfikacja dna morskiego}
\author{Igor Józefowicz, indeks 193257}
\date{}

\begin{document}
\maketitle

\section*{1. Dane i cel}
Raport dotyczy ćwiczenia z przedmiotu \emph{Uczenie maszynowe w badaniach Ziemi}. Celem było obliczenie parametrów opisujących sondowania BSS z sonaru wielowiązkowego, pokazanie ich rozkładów oraz sprawdzenie skuteczności prostych klasyfikatorów dla czterech typów dna.

Wykorzystano plik \texttt{dane\_bss.mat}, który zawiera tablicę \texttt{vbss} o rozmiarze \(4 \times 600 \times 160\): cztery typy dna, 600 sondowań dla każdego typu i 160 wiązek sonaru w jednym sondowaniu. Do obliczeń użyto pierwszych 500 sondowań z każdej klasy, zgodnie z dostarczonymi skryptami.

Typy dna: 1 -- muł, 2 -- mieszanina mułu, piasku i odpadów antropogenicznych, 3 -- piasek drobnoziarnisty, 4 -- piasek gruboziarnisty.

Najlepszy wynik klasyfikacji w wykonanych próbach uzyskał model \texttt{knn\_5} dla wszystkich parametrów i 60\% danych uczących: accuracy=96.00\%.

\section*{2. Obliczone parametry}
Obliczono cztery parametry wymagane w instrukcji oraz trzy lekkie rozszerzenia. W tabeli zastosowano skrócone nagłówki: GLCM E oznacza entropię, GLCM LH spójność lokalną, L-R różnicę sektorów lewego i prawego.

\begin{table}[H]
\centering
\scriptsize
\begin{tabular}{lrrrrrrr}
\toprule
typ dna & mean & drop & GLCM E & GLCM LH & std & L-R & edge-center \\
\midrule
1. muł & -33.29 & 7.17 & 0.58 & 0.85 & 2.71 & 3.59 & -15.28 \\
2. muł/piasek/odpady & -25.16 & 0.89 & 0.52 & 0.86 & 3.38 & 0.14 & -3.92 \\
3. piasek drobnoziarnisty & -23.68 & 6.63 & 0.77 & 0.77 & 4.26 & 1.99 & -6.76 \\
4. piasek gruboziarnisty & -24.34 & 9.64 & 0.70 & 0.77 & 4.58 & 3.25 & -5.70 \\
\bottomrule
\end{tabular}
\end{table}

\section*{3. Wizualizacja parametrów}
Wykresy pokazują, że klasa 1 jest najlepiej oddzielona od pozostałych przez poziom BSS. Klasy 2, 3 i 4 częściowo nakładają się w przestrzeni parametrów, dlatego sama średnia BSS nie wystarcza do stabilnej klasyfikacji. Dodatkowe parametry opisują kształt zmian sygnału wraz z kątem sondowania i lokalną teksturę obrazu sonarowego.


\begin{figure}[H]
\centering
\includegraphics[width=0.86\textwidth]{\detokenize{wyniki_lab6/01_bss_images.png}}
\caption{Obrazy BSS dla czterech typów dna.}
\end{figure}


\begin{figure}[H]
\centering
\includegraphics[width=0.78\textwidth]{\detokenize{wyniki_lab6/02_scatter_mean_homogeneity.png}}
\caption{Rozkład 2D: średnia BSS i spójność lokalna GLCM.}
\end{figure}


\begin{figure}[H]
\centering
\includegraphics[width=0.78\textwidth]{\detokenize{wyniki_lab6/03_scatter_mean_edge_center.png}}
\caption{Rozkład 2D: średnia BSS i różnica sektorów.}
\end{figure}


\begin{figure}[H]
\centering
\includegraphics[width=0.92\textwidth]{\detokenize{wyniki_lab6/04_boxplots.png}}
\caption{Rozkłady wybranych parametrów według typu dna.}
\end{figure}


\begin{figure}[H]
\centering
\includegraphics[width=0.78\textwidth]{\detokenize{wyniki_lab6/05_correlation_heatmap.png}}
\caption{Korelacje między obliczonymi parametrami.}
\end{figure}


\begin{figure}[H]
\centering
\includegraphics[width=0.82\textwidth]{\detokenize{wyniki_lab6/06_scatter_3d.png}}
\caption{Dodatkowy rozkład 3D dla trzech parametrów.}
\end{figure}


\section*{4. Klasyfikacja}
Klasyfikację minimalnoodległościową odtworzono przez \texttt{NearestCentroid}. Dane normalizowano metodą min-max, a podział uczący/testowy wykonano per klasa: pierwsze 20\%, 40\% albo 60\% obiektów jako zbiór uczący, reszta jako testowy.

\subsection*{Pary parametrów}
\begin{table}[H]
\centering
\small
\begin{tabular}{llrr}
\toprule
model & parametry & trening [\%] & accuracy [\%] \\
\midrule
min. odl. & mean + GLCM LH & 40 & 77.58 \\
min. odl. & mean + GLCM LH & 60 & 75.25 \\
min. odl. & mean + drop & 60 & 70.88 \\
min. odl. & mean + edge-center & 20 & 69.62 \\
min. odl. & mean + drop & 40 & 68.92 \\
min. odl. & mean + drop & 20 & 67.44 \\
min. odl. & mean + edge-center & 40 & 67.17 \\
min. odl. & mean + edge-center & 60 & 65.12 \\
min. odl. & mean + GLCM LH & 20 & 63.19 \\
\bottomrule
\end{tabular}
\end{table}

\subsection*{Wszystkie parametry}
\begin{table}[H]
\centering
\small
\begin{tabular}{llrr}
\toprule
model & parametry & trening [\%] & accuracy [\%] \\
\midrule
kNN & wszystkie & 60 & 96 \\
kNN & wszystkie & 40 & 92.50 \\
min. odl. & wszystkie & 60 & 89.75 \\
kNN & wszystkie & 20 & 86.69 \\
min. odl. & wszystkie & 40 & 84.75 \\
min. odl. & wszystkie & 20 & 76.44 \\
\bottomrule
\end{tabular}
\end{table}

Najlepsza para dla klasyfikatora minimalnoodległościowego przy 20\% treningu to \texttt{mean+edge\_center}, accuracy=69.62\%.


\begin{table}[H]
\centering
\small
\caption{Najlepsza para parametrów, 20\% treningu}
\begin{minipage}{0.47\textwidth}
\centering
\textbf{Liczba obiektów}\\[0.35em]
\begin{tabular}{lrrrr}
\toprule
 & kl. 1 & kl. 2 & kl. 3 & kl. 4 \\
\midrule
rzecz. 1 & 400 & 0 & 0 & 0 \\
rzecz. 2 & 0 & 245 & 21 & 134 \\
rzecz. 3 & 0 & 18 & 328 & 54 \\
rzecz. 4 & 0 & 68 & 191 & 141 \\
\bottomrule
\end{tabular}
\end{minipage}\hfill
\begin{minipage}{0.47\textwidth}
\centering
\textbf{Udział procentowy [\%]}\\[0.35em]
\begin{tabular}{lrrrr}
\toprule
 & kl. 1 & kl. 2 & kl. 3 & kl. 4 \\
\midrule
rzecz. 1 & 100.00 & 0.00 & 0.00 & 0.00 \\
rzecz. 2 & 0.00 & 61.25 & 5.25 & 33.50 \\
rzecz. 3 & 0.00 & 4.50 & 82.00 & 13.50 \\
rzecz. 4 & 0.00 & 17.00 & 47.75 & 35.25 \\
\bottomrule
\end{tabular}
\end{minipage}
\end{table}


\begin{table}[H]
\centering
\small
\caption{NearestCentroid, wszystkie parametry, 20\% treningu}
\begin{minipage}{0.47\textwidth}
\centering
\textbf{Liczba obiektów}\\[0.35em]
\begin{tabular}{lrrrr}
\toprule
 & kl. 1 & kl. 2 & kl. 3 & kl. 4 \\
\midrule
rzecz. 1 & 400 & 0 & 0 & 0 \\
rzecz. 2 & 0 & 271 & 40 & 89 \\
rzecz. 3 & 0 & 0 & 333 & 67 \\
rzecz. 4 & 0 & 0 & 181 & 219 \\
\bottomrule
\end{tabular}
\end{minipage}\hfill
\begin{minipage}{0.47\textwidth}
\centering
\textbf{Udział procentowy [\%]}\\[0.35em]
\begin{tabular}{lrrrr}
\toprule
 & kl. 1 & kl. 2 & kl. 3 & kl. 4 \\
\midrule
rzecz. 1 & 100.00 & 0.00 & 0.00 & 0.00 \\
rzecz. 2 & 0.00 & 67.75 & 10.00 & 22.25 \\
rzecz. 3 & 0.00 & 0.00 & 83.25 & 16.75 \\
rzecz. 4 & 0.00 & 0.00 & 45.25 & 54.75 \\
\bottomrule
\end{tabular}
\end{minipage}
\end{table}


\begin{table}[H]
\centering
\small
\caption{kNN, wszystkie parametry, 20\% treningu}
\begin{minipage}{0.47\textwidth}
\centering
\textbf{Liczba obiektów}\\[0.35em]
\begin{tabular}{lrrrr}
\toprule
 & kl. 1 & kl. 2 & kl. 3 & kl. 4 \\
\midrule
rzecz. 1 & 400 & 0 & 0 & 0 \\
rzecz. 2 & 0 & 323 & 36 & 41 \\
rzecz. 3 & 0 & 1 & 344 & 55 \\
rzecz. 4 & 0 & 0 & 80 & 320 \\
\bottomrule
\end{tabular}
\end{minipage}\hfill
\begin{minipage}{0.47\textwidth}
\centering
\textbf{Udział procentowy [\%]}\\[0.35em]
\begin{tabular}{lrrrr}
\toprule
 & kl. 1 & kl. 2 & kl. 3 & kl. 4 \\
\midrule
rzecz. 1 & 100.00 & 0.00 & 0.00 & 0.00 \\
rzecz. 2 & 0.00 & 80.75 & 9.00 & 10.25 \\
rzecz. 3 & 0.00 & 0.25 & 86.00 & 13.75 \\
rzecz. 4 & 0.00 & 0.00 & 20.00 & 80.00 \\
\bottomrule
\end{tabular}
\end{minipage}
\end{table}


\section*{5. Dodany kod własny}
Pliki startowe do laboratorium zawierały skrypty MATLAB/Octave, m.in. \texttt{vbss.m}, \texttt{obl\_srednia.m}, \texttt{obl\_spadek\_dB.m}, \texttt{obl\_GLCM\_entr\_spoj.m}, \texttt{rozkl2D.m} oraz \texttt{kl\_min\_dist\_2D.m}. Nowy kod dodany w ramach rozwiązania znajduje się w pliku \texttt{lab6\_solution.py}. Poniżej pokazano tylko najważniejsze fragmenty: obliczenie parametrów, obliczenie parametrów GLCM oraz wspólną funkcję klasyfikacji.


\begin{lstlisting}[caption={Fragment własnego kodu: obliczanie parametrów BSS.}]
def compute_features(vbss):
    values = replace_zero_with_min(vbss)
    arrays = {
        "mean_80_100": values[:, :N_OBJECTS, beam_slice(80, 100)].mean(axis=2),
        "drop_80_100": values[:, :N_OBJECTS, 79] - values[:, :N_OBJECTS, 99],
        "std_80_100": values[:, :N_OBJECTS, beam_slice(80, 100)].std(axis=2),
        "left_right_diff": values[:, :N_OBJECTS, beam_slice(61, 80)].mean(axis=2)
            - values[:, :N_OBJECTS, beam_slice(91, 100)].mean(axis=2),
        "edge_center_diff": values[:, :N_OBJECTS, beam_slice(1, 20)].mean(axis=2)
            - values[:, :N_OBJECTS, beam_slice(80, 100)].mean(axis=2),
    }
\end{lstlisting}



\begin{lstlisting}[caption={Fragment własnego kodu: obliczanie entropii i spójności lokalnej GLCM.}]
for class_idx in range(4):
    for object_idx in range(N_OBJECTS):
        window = values[class_idx, object_idx : object_idx + 30, beam_slice(91, 100)]
        entropy, homogeneity = glcm_entropy_homogeneity(
            window, levels=8, distance=1, total_min=total_min, total_max=total_max
        )
        glcm_entropy[class_idx, object_idx] = entropy
        glcm_homogeneity[class_idx, object_idx] = homogeneity
\end{lstlisting}



\begin{lstlisting}[caption={Fragment własnego kodu: wspólna funkcja klasyfikacji i ewaluacji.}]
def run_classifier(df, feature_names, train_pct, classifier_name):
    x = MinMaxScaler().fit_transform(df[feature_names].to_numpy())
    y = df["class_id"].to_numpy()
    train_idx, test_idx = split_by_class(train_pct, y)
    classifier = NearestCentroid() if classifier_name == "nearest_centroid" else KNeighborsClassifier(n_neighbors=5)
    classifier.fit(x[train_idx], y[train_idx])
    predictions = classifier.predict(x[test_idx])
    return accuracy_score(y[test_idx], predictions) * 100, confusion_matrix(y[test_idx], predictions)
\end{lstlisting}


\section*{6. Wnioski}
\begin{enumerate}
    \item Klasa 1, czyli muł, jest najłatwiejsza do rozpoznania. W głównych macierzach niezgodności wszystkie lub prawie wszystkie obiekty tej klasy trafiają do poprawnej klasy.
    \item Największe pomyłki występują między klasami 2, 3 i 4, co zgadza się z wykresami parametrów, gdzie rozkłady tych klas częściowo się nakładają.
    \item Najlepsza para parametrów przy 20\% treningu uzyskała accuracy=69.62\%, więc klasyfikacja 2D jest użyteczna, ale ograniczona.
    \item Najlepszy wynik całego zestawienia to \texttt{knn\_5} dla wszystkich parametrów, train\_pct=60\%, accuracy=96.00\%.
    \item Dla \texttt{NearestCentroid} na wszystkich parametrach accuracy wzrosło od 76.44\% do 89.75\%, więc większy zbiór uczący poprawił stabilność klasyfikacji.
    \item Dodatkowe parametry własne są przydatne, bo opisują nie tylko średni poziom BSS, lecz także zmianę sygnału wraz z kątem sondowania i lokalną teksturę obrazu.
\end{enumerate}

\end{document}
